\documentclass{beamer}
\usepackage[utf8]{inputenc}
\usetheme{Boadilla}      % 主题
\usecolortheme{default} % 颜色主题

\usepackage{mathspec}
\usepackage{amssymb}
\usepackage{xcolor}
\usepackage{graphicx}
\usepackage{float}
\usepackage{booktabs}
\usepackage{tabularx}
\usepackage{caption}
\usepackage{multirow}
\usepackage{array}
\usepackage{colortbl}

\usepackage{tikz}
\tikzset{x=10mm, y=10mm}
\usetikzlibrary{arrows.meta, positioning, fit, backgrounds}

\usepackage{listings}
% 自定义列表环境
\lstset{
    numbers=left,           % 行号在左侧
    % numbersep=5pt,          % 行号与代码的间距
    %xleftmargin=0pt,        % 左边距 (行号从这里开始算) 
    % xrightmargin=0pt,
    frame=single,           % 添加边框 (可选) 
    % framexleftmargin=0pt,
    breaklines=true,
    captionpos=b,           % 关键：将 caption 放在底部 (b = bottom) 
    % aboveskip=10pt,         % 代码块上方间距
    % belowskip=10pt,         % 代码块下方间距
}

\usepackage{makecell}
\usepackage{setspace}
\setstretch{1}

\usepackage{fontspec}
\usepackage{xeCJK}
\setmainfont{Times New Roman}
\setsansfont{Arial}
\setmonofont{Courier New}
\setCJKmainfont{Songti SC}[
    AutoFakeBold=true,
    AutoFakeSlant=true
]
\setCJKsansfont{PingFang SC}[
    AutoFakeBold=true,
    AutoFakeSlant=true
]
\setCJKmonofont{Songti SC}[
    AutoFakeBold=true,
    AutoFakeSlant=true
]


\usepackage{hyperref}

% 设置正文大小（相当于全局基准）
\setbeamerfont{normal text}{size=\small}   % 将正文设为 \small 大小

% 设置各级标题大小
\setbeamerfont{frametitle}{size=\Large}    % 每页标题用 \Large
\setbeamerfont{title}{size=\huge}          % 首页大标题用 \huge
\setbeamerfont{author}{size=\normalsize}   % 首页作者名用正常大小

% 设置目录页字体大小
\setbeamerfont{section in toc}{size=\normalsize}    % 一级目录用正常大小
\setbeamerfont{subsection in toc}{size=\footnotesize} % 二级目录用小字

% 为了让 \setbeamerfont{normal text} 生效，必须加上这句
\renewcommand{\normalsize}{\usebeamerfont{normal text}\selectfont}

% 目录风格
\setbeamertemplate{section in toc}{
  \leavevmode
  \leftskip=0pt
  (\inserttocsectionnumber).\ % 这里改成你想要的格式
  \inserttocsection
  \par
}

% 自定义顶栏
\setbeamertemplate{headline}{
  \leavevmode
  \begin{beamercolorbox}[wd=\paperwidth,ht=2.5ex,dp=1.125ex]{section in head/foot}
    \hspace{1em}%
    \insertsectionhead
    \hfill
    \insertshorttitle  % 显示短标题，不会报错
    \hspace*{1em}%
  \end{beamercolorbox}
}

% 标题信息
\title{声指令随 - EpochPoint}
\subtitle{北京邮电大学智能机器人实验期末大作业}
\author[BUPT]{
  任志诚 2023212020
  梁恒 2023123456
  金泽林 2023212012
}
\institute[]{BUPT}
\date{\today}

\begin{document}

    % 标题页
    \begin{frame}
        \titlepage
    \end{frame}

    % 目录页
    \begin{frame}{Content}
        \tableofcontents
    \end{frame}

    % 第一部分
    \section{基于LangChain架构的智能体}

    \begin{frame}{Langchain是什么？}
        LangChain 是一个用于构建大语言模型（LLM）应用的主流开源框架，核心是连接模型、数据与工具，解决 LLM 孤立、无记忆、难执行复杂任务的问题。总而言之，可以用起开发agent从而实现 (在我们的设计中) 
        如下功能
        \begin{center}
        \resizebox{\textwidth}{!}{
        \begin{tikzpicture}[
            >=Stealth,
            % 样式定义
            host/.style={
                rectangle, draw=blue!70!black, rounded corners, 
                minimum height=2.2cm, minimum width=2.8cm, 
                fill=blue!5, align=center,
                font=\sffamily\bfseries
            },
            cloud/.style={
                rectangle, draw=orange!70!black, rounded corners,
                minimum height=2.2cm, minimum width=2.8cm,
                fill=orange!5, align=center,
                font=\sffamily\bfseries
            },
            process/.style={
                rectangle, draw=gray!60, dashed, rounded corners=2pt,
                minimum height=0.7cm, minimum width=2.8cm,
                fill=white, align=center,
                font=\sffamily\footnotesize
            },
            processgroup/.style={
                rectangle, draw=gray!40, rounded corners,
                inner sep=6pt, fill=gray!3, align=center
            }
        ]

        % ========== 左右主体 ==========
        \node[host] (host) {本地 Host\\ (Agent)};
        \node[cloud, right=4cm of host] (llm) {云端 LLM};

        % ========== 云端内部流程（放在 LLM 右侧）==========
        % 流程组框
        \node[processgroup, right=1.2cm of llm] (processbox) {
            \begin{tikzpicture}
                \node[process] (p1) {图像接收};
                \node[process, below=0.15cm of p1] (p2) {图文特征融合};
                \node[process, below=0.15cm of p2] (p3) {LLM 语义理解};
                \node[process, below=0.15cm of p3] (p4) {生成图像描述};
                
                % 流程连线
                \draw[->, thick, gray!60] (p1) -- (p2);
                \draw[->, thick, gray!60] (p2) -- (p3);
                \draw[->, thick, gray!60] (p3) -- (p4);
            \end{tikzpicture}
        };

        % 流程组标题
        \node[above=0.3cm of processbox.north, font=\sffamily\small\color{gray!70!black}] 
            {内部处理流程};

        % 发送：Host -> LLM
        \draw[->, thick, blue!70!black] 
            (host.east) -- 
            node[above, font=\sffamily\small] {发送 Image + Prompts} 
            (llm.west);

        % 返回：LLM -> Host
        \draw[<-, thick, red!70!black] 
            (host.east) -- 
            node[below, font=\sffamily\small] {返回图像解析内容} 
            (llm.west);

        % LLM 到流程组的交互箭头
        \draw[->, thick, orange!70!black, dashed] 
            (llm.east) -- 
            node[above, font=\sffamily\small] {处理} 
            (processbox.west);

        \end{tikzpicture}
        }
    \end{center}
    \end{frame}

    \begin{frame}{使用LangChain的必要性}
        调研机器人配置，得到如下信息：
        \begin{enumerate}
            \item[(1).] CPU：$4$核$8$线程 (i5-8259U)，无独立显卡，本地推理大模型能力受限。
            \item[(2).] 内存与存储：仅$8\text{GB}$内存，$116\text{GB}$系统盘，难以加载或运行$7\text{B}$以上量级的LLM。
            \item[(3).] 系统环境：Ubuntu 16.04（较老），Python 2.7/3.5为主，不直接支持多数新LLM框架。
            \item[(4).] 开机后负载极低（CPU空闲$> 0.97$），主要运行MQTT消息服务器(EMQX)，适合作为调用云端LLM的轻量Agent。
        \end{enumerate}

        可见，机器人的配置十分低端，根本无法进行本地图像识别；为了配置本地和云端LLM之间的通信，利用LangChain搭建一个简单的agent是一个可行的方案

    \end{frame}

    \begin{frame}{具体的实现方案}
        由于机器人全局环境为python3.8，而LangChain推荐的python环境为3.10$+$。故这里在Docker中创建相关环境
        
    \end{frame}

\end{document}